% \iffalse meta-comment
%
% hit-bib.dtx 是参考文献页
%
% \fi
%
% \section{参考文献页}
%
% ^^A ======================================================================
% ^^A Module thesis-bib: 参考文献加载与 natbib/gbt7714 设置（hit-thesis）
% ^^A ======================================================================
%    \begin{macrocode}
%<*thesis-bib>
\RequirePackage[sort&compress]{gbt7714}
\RequirePackage{etoolbox}
\RequirePackage{xparse}
\cs_new:Npn \bibstyle@numerical
  { \bibpunct {[} {]} {,} {s} {,} { \textsuperscript {,} } }
%    \end{macrocode}
%
% 删除inline等定义，此命令与 \pkg{gbt7714} 冲突
% \changes{v3.0.15}{2021/05/06}{删除inline等定义，此命令与 \pkg{gbt7714} 冲突}
%    \begin{macrocode}
% \newcommand\bibstyle@authoryear{\bibpunct{(}{)}{;}{a}{,}{,}}
%    \end{macrocode}
%
% \changes{v2.0.06}{2018/12/05}{在\cs{inlinecite}内添加空格}
%    \begin{macrocode}
% \newcommand\bibstyle@inline{\bibpunct{[}{]}{,}{n}{,}{\hit@inline@separator}}
\citestyle{numerical}
\DeclareRobustCommand\inlinecite{\@inlinecite}
\cs_set:Npn \@inlinecite #1
  {
    \group_begin:
      \citestyle { inline }
      \cs_set_eq:NN \@cite \NAT@citenum
      \citep {#1}
    \group_end:
  }
\cs_set_eq:NN \onlinecite \inlinecite
%    \end{macrocode}
%
% \changes{v3.1a}{2022/05/17}{按照 \hitszthesis{} 生成深圳校区的参考文献样式 \texttt{hitszthesis.bst}。\issue[97]}
%    \begin{macro}{\@biblabel}
% 设置参考文献列表的序号中的括号为中文括号。
% \changes{v3.1b}{2022/06/06}{将 ［］前加上 \cs{mbox\{\}} 来获得正确的序号长度。} 
% 在［\#1］前加上 \cs{mbox\{\}} 来获得正确的序号长度，否则会因为 \XeLaTeX 下的标点挤压而计算错误，这里应该整个包裹起来，否则换行后缩进会有问题。
% \changes{v3.1d}{2025/03/03}{按着 \issue[211] 来修改，将 \cs{mbox\{\}} 包住整个 label。}
% \changes{v3.1d}{2025/03/03}{按照 \issue[215] 将哈尔滨硕博毕业论文的参考文献序号也改为全角。}
%    \begin{macrocode}
\bool_lazy_or:nnT
  { \bool_if_p:N \g__hit_shenzhen_bool }
  {
    \bool_lazy_and_p:nn
      { \bool_if_p:N \g__hit_harbin_bool }
      {
        \bool_lazy_or_p:nn
          { \bool_if_p:N \g__hit_master_bool } { \bool_if_p:N \g__hit_doctor_bool }
      }
  }
  { \cs_set:Npn \@biblabel #1 { \mbox { ［#1］ } } }
%    \end{macrocode}
% \end{macro}
% 
% \begin{environment}{thebibliography}
% 修改默认的 \env{thebibliography} 环境来适应间距的要求。
% \changes{v3.2a}{0000/00/00}{设置参考文献列表标号的对齐方式}
%    \begin{macrocode}
\RenewDocumentEnvironment{thebibliography}{m}{%
%    \end{macrocode}
%
% \changes{v3.1d}{2025/03/03}{更改完全角符号后的半角点号可能出现在行首，规避一下。}
% 这里我直接贴一个 issue 的链接：\href{https://github.com/CTeX-org/ctex-kit/issues/701}{CTeX-Kit \#701}
% \changes{v3.2a}{0000/00/00}{FullRight引入了太多bug, 取消之 (bst文件中加入了nolinebreak)}
%    \begin{macrocode}
  % \xeCJKDeclareCharClass{FullRight}{"002E}%
  \hit@clear@page\phantomsection%
  \hit@appendix@chapter*{\bibname}[\hit@bibliography@heading@en]
  \normalsize
  \list{\@biblabel{\@arabic\c@enumiv}}{%
    \str_if_eq:VnT \g_bib_labelalign_str {left} {
      \RenewDocumentCommand\makelabel{m}{##1\hfill}
    }
    \settowidth{\labelwidth}{\@biblabel{#1}}
%    \end{macrocode}
% \changes{v3.1b}{2022/06/06}{修改深圳校区的 \cs{labelsep} 为 0pt。}
%    \begin{macrocode}
    \hit@if@option@TF{shenzhen}{\dim_set:Nn \labelsep {0pt}}{\dim_set:Nn \labelsep {0.5em}}
    \dim_set:Nn \itemindent {0pt}
    \dim_set:Nn \leftmargin {\labelsep+\labelwidth}
    \dim_add:Nn \itemsep {-0.8em}
    \usecounter{enumiv}%
    \cs_set_eq:NN \p@enumiv \@empty
  \cs_set:Npn \theenumiv {\@arabic\c@enumiv}}%
  \sloppy\frenchspacing
%    \end{macrocode}
% \changes{v2.0.07}{2019/03/02}{添加flushbottom到thebibliography环境中}
%    \begin{macrocode}
  \flushbottom
%    \end{macrocode}
% \changes{v2.0.03}{2018/10/08}{添加参考文献分割开关}
%    \begin{macrocode}
  \hit@if@option@TF{splitbibitem}%
  {%
  \clubpenalty0
  \@clubpenalty \clubpenalty
  \widowpenalty0%
  \interlinepenalty-50%
  }%
  {%
  \clubpenalty4000
  \@clubpenalty \clubpenalty
  \widowpenalty4000%
  \interlinepenalty4000%
  }%
\sfcode`\.\@m}
{\cs_set_nopar:Npn \@noitemerr
  {\@latex@warning{Empty `thebibliography' environment}}%
\endlist\frenchspacing}
%    \end{macrocode}
% \end{environment}
%    \begin{macrocode}
\cs_set:Npn \NAT@cite #1#2#3
      {\ifNAT@swa\NAT@@open\if*#2*\else#2\NAT@spacechar\fi
        #1\NAT@@close\if*#3*\else\textsuperscript{#3}\fi\else#1\fi\endgroup}
%    \end{macrocode}
% \changes{v3.2a}{0000/00/00}{控制参考文献列表中作者是否全大写和最大显示数量}
%    \begin{macrocode}
%% 用以控制参考文献列表中单个作者的格式
\exp_args:NNe \cs_new:Nn \g_bib_author_cs:n {
    \bool_if:NTF \g_bib_authorupper_bool {
        \exp_not:N \text_uppercase:n {#1}
    } {
        {#1}
    }
}
%% 用以输出参考文献列表中的作者
\NewDocumentCommand{\bibauthor}{m}{
  \int_zero:N \l_tmpa_int
  \exp_args:Ne \tl_map_inline:nn {
    \exp_args:Ne \tl_tail:n {#1}
  } {
    \int_compare:nNnT
      \l_tmpa_int = \g_bib_maxauthor_int {
      \prg_break:n {\tl_head:N {#1}}
    }
    \int_incr:N \l_tmpa_int
    \g_bib_author_cs:n {##1}
  }
  \prg_break_point:
}
%    \end{macrocode}
%
% \changes{v3.2a}{0000/00/00}{\cs{citeup}（上标引用宏）从 book-floats 迁入 book-bib，归位至引用语义所属模块}
% 上标引用宏 \cs{citeup}（语义属引用,非浮动体）。
%    \begin{macrocode}
\NewDocumentCommand \citeup { m } { \textsuperscript { \cite {#1} } }
%</thesis-bib>
%    \end{macrocode}
%
% \Finale
